\section{Introduction}

An arc in a projective plane is a set of points no three of which are
collinear.  Ordinary completeness asks whether one further point can be
adjoined.  If the arc is preserved by a field automorphism, the natural
equivariant continuation unit is instead a complete Galois orbit.  Over a
quadratic extension this means either a fixed point or a nonfixed conjugate
pair.

Let $F=\F_s$, let $K=\F_{s^2}$, and let
\[
  \phi\colon [x:y:z]\longmapsto[x^s:y^s:z^s]
\]
act on $\PG(2,K)$.  The fixed points and fixed lines form the embedded Baer
subplane $\PG(2,F)$.  This paper studies the following question: when must a
$\phi$-invariant arc admit an extension by a fresh pair
$\{P,\phi(P)\}$?

Under the arc--code correspondence, a $k$-arc gives the projective columns
of a $[k,3,k-2]_{s^2}$ MDS code.  Adjoining $\{P,\phi(P)\}$ is therefore a
Frobenius-compatible two-column MDS extension, and alternate-orbit repair is
replacement of an erased conjugate column pair.  Here ``repair'' changes the
generator-column configuration; it is not erasure decoding inside a fixed
code.

\begin{theorem}[Robust extension of invariant eight-arcs]
\label{thm:main}
Let $s\ge5$ be a prime power.  Every $\phi$-invariant eight-arc in
$\PG(2,s^2)$ admits a fresh nonfixed conjugate-pair extension.  More
precisely:
\begin{enumerate}[label=(\roman*),leftmargin=2.1em]
\item over $\F_{25}$ there are at least four such pairs, so deleting any
      selected nonfixed orbit from an invariant ten-arc leaves at least
      three alternate repairs;
\item for every $s\ge7$ there are at least $319$ legal pairs, so deletion of
      any selected nonfixed orbit from an invariant ten-arc leaves at least
      $318$ alternate repairs.
\end{enumerate}
\end{theorem}

The proof has one organizing mechanism.  A nonfixed pair lies on a unique
fixed mate line, and a legal pair must lie on a fixed line disjoint from the
old arc.  Each noninvariant old-secant orbit forbids at most one candidate on
such a carrier.  Counting empty carriers gives the uniform criterion.  The
exact correction records the two ways in which the first-order count
overestimates the obstruction: the secant orbit can be invisible because
its fixed center lies on the carrier, or several visible orbits can charge
the same pair.

The argument first proves the carrier count and its exact correction.  For
$s\ge7$ those structural bounds imply Theorem~\ref{thm:main}(ii).  At $s=5$
one profile lies beyond the first-order range; two projective normalizations
reduce it to a finite residual problem.  Proposition~\ref{prop:f2} gives the
exact minimum $32$ in that profile and classifies its normalized equality
cases.  Combining it with the four structural profiles proves
Theorem~\ref{thm:main}(i).

The quantitative criterion and its exact correction are the structural
part of the paper.  Proposition~\ref{prop:f2} is the only exhaustive finite
step.  Its lower bound, attainment, and equality-orbit exhaustion are kernel
checked in the normalized coordinate model; the manuscript carries those
facts through the two projective normalizations.  Section~\ref{sec:formalization}
states the certificate and trust boundary.  The parameterized exchange
theorem identifies the general phase
$\lfloor(k-1)^2/4\rfloor+r+1\le s(s-1)/2$.  The saturation corollary and the
Clebsch specialization are secondary consequences of the same carrier
count.  The saturation scale is the classical Lunelli--Sce $\sqrt{2q}$
scale~\cite{LunelliSce1964,BallLavrauw2019}; no new constant or sharpness
claim is made.

The classical inputs are point and line counts in a projective plane, the
fixed Baer subplane, Hilbert--90 normalization, and secant counting; see
Hirschfeld~\cite{Hirschfeld1998} and Martin~\cite{Martin1967}.

\subsection*{Prior art and claim boundary}

Baker and Wantz~\cite{BakerWantz2005} use Frobenius-invariant arcs in the
Hughes plane and, in a maximality argument, consider a point together with
its Frobenius image.  This is genuine precedence for the conjugate-point
maneuver.  Their result does not count legal conjugate pairs on empty Baer
lines.  Jungnickel's arcs in Baer-semiplane complements
\cite{Jungnickel1987} and Ueberberg's geometries of Frobenius collineations
\cite{Ueberberg1997} are also adjacent but concern different incidence or
collineation problems.

Dye gives exact completion geometry for the special family of Clebsch
hexagons and a five-valent adjacency graph defined by sharing an associated
triangle~\cite[Proposition~1 and Theorems~7--8, pp.~282--284]{Dye1991}.
Blokhuis, Seress, and Wilbrink classify extremal complete exterior sets
relative to a conic~\cite[pp.~143--146]{BlokhuisSeressWilbrink1992}.
Neither work treats arbitrary plane arcs invariant under the quadratic
field-Frobenius involution, counts fresh legal conjugate-pair extensions, or
requires a different replacement after every selected-orbit deletion.

Ordinary arc and MDS lengthening results
\cite{Alderson2005,AldersonBruenSilverman2007,Alderson2026} impose no
quadratic Galois-pair constraint.  Likewise, the classification of
Galois-invariant subsets of $\mathbf P^1$~\cite{LopezMaisnerNartXarles2002}
does not treat arbitrary plane arcs or their pair extensions.  Bounded
zbMATH Open, OpenAlex, Crossref, and source-level searches found no exact
precursor for the quantitative criterion below for arbitrary quadratic-
field-Frobenius-invariant arcs, or for the uniform $\PG(2,25)$ theorem.
This is negative search evidence, not a claim of historical priority.
